% !PRINTED 183
\weberpage{183}
Man erhält dann aus den Entwickelungen in \S\,25 mit Rücksicht darauf, daß nach \S\,34 und \S\,42 für ein verschwindendes $q$
\[
\eta(\omega) = q^{\frac{1}{12}}, \qquad 2K = \pi
\]
wird:
\weq{5}{\begin{array}{r@{}l@{\;}l}
e^{\frac{v^2}{6}} & \sin v & = \displaystyle\sum_{0,\infty}^{\nu} A_\nu(0)\,\dfrac{v^{2\nu+1}}{(2\nu+1)!},\\[8pt]
e^{\frac{v^2}{6}} & & = \displaystyle\sum_{0,\infty}^{\nu} B_\nu(0)\,\dfrac{v^{2\nu}}{(2\nu)!},\\[8pt]
e^{\frac{v^2}{6}} & & = \displaystyle\sum_{0,\infty}^{\nu} C_\nu(0)\,\dfrac{v^{2\nu}}{(2\nu)!},\\[8pt]
e^{\frac{v^2}{6}} & \cos v & = \displaystyle\sum_{0,\infty}^{\nu} D_\nu(0)\,\dfrac{v^{2\nu}}{(2\nu)!}.
\end{array}}
Die Formeln (2) ergeben nun:
\weq{6}{\begin{array}{l}
A_\nu(\varkappa'^2) = (-1)^\nu A_\nu(\varkappa^2),\\[4pt]
B_\nu(\varkappa'^2) = (-1)^\nu B_\nu(\varkappa^2),\\[4pt]
C_\nu(\varkappa'^2) = (-1)^\nu D_\nu(\varkappa^2),\\[4pt]
D_\nu(\varkappa'^2) = (-1)^\nu C_\nu(\varkappa^2),
\end{array}}
woraus folgt, daß auch für $\varkappa^2 = 1$ diese Funktionen endlich bleiben, nämlich:
\weq{7}{\begin{array}{l}
A_\nu(1) = (-1)^\nu A_\nu(0),\\[4pt]
B_\nu(1) = (-1)^\nu B_\nu(0),\\[4pt]
C_\nu(1) = (-1)^\nu D_\nu(0),\\[4pt]
D_\nu(1) = (-1)^\nu C_\nu(0).
\end{array}}
Das System (3) läßt sich ferner so schreiben:
\weq{8}{\begin{array}{l}
\varkappa'^{2\nu} A_\nu\left(\dfrac{-\varkappa^2}{\varkappa'^2}\right) = A_\nu(\varkappa^2),\\[6pt]
\varkappa'^{2\nu} B_\nu\left(\dfrac{-\varkappa^2}{\varkappa'^2}\right) = C_\nu(\varkappa^2),\\[6pt]
\varkappa'^{2\nu} C_\nu\left(\dfrac{-\varkappa^2}{\varkappa'^2}\right) = B_\nu(\varkappa^2),\\[6pt]
\varkappa'^{2\nu} D_\nu\left(\dfrac{-\varkappa^2}{\varkappa'^2}\right) = D_\nu(\varkappa^2),
\end{array}}
woraus für ein unendliches $\varkappa^2$:
\weq{9}{\begin{array}{l}
A_\nu(\varkappa^2) = \varkappa^{2\nu} A_\nu(0),\\[4pt]
B_\nu(\varkappa^2) = \varkappa^{2\nu} D_\nu(0),\\[4pt]
C_\nu(\varkappa^2) = \varkappa^{2\nu} B_\nu(0),\\[4pt]
D_\nu(\varkappa^2) = \varkappa^{2\nu} C_\nu(0),
\end{array}}
%CONT>>
